\subsection{社区原始CLI与补丁的隔离执行}
选择D1中先前已审查的三个包：数据质量审计、资深数据科学及资深机器学习工程，合计九个原始入口。
它们来自同一仓库，部分共享模板，故不是九个独立自然Skill样本。
其中数据质量包的三个原始脚本完整按冻结版本复制，另外六个入口用于核对声明完成与实际输出之间的关系。
不向上游提交修改，也不覆盖原始1,881对象交付。

CSV试验沿用四份UCI数据，四个条件（普通、引号内逗号/换行、字面NA、纯空格）与1/8/48行形成48个输入；
另加一个可手算的阈值回归例，共49个不同输入。
画像、画像修复、缺失分析各运行一次；异常检测的原始与修复版本各在1.5、3.0阈值运行一次。
再对六个模板入口分别传入有效及缺失输入，共355次CLI进程执行。
补丁只改变单样本标准差处理与IQR阈值传递各一处表达式，修改前后都有源码指纹。

所有社区代码只在公开CI的固定Docker镜像内运行：无网络、非root、根文件系统与输入只读、去能力、
禁止新增权限，128 MiB、1 CPU、32 PID，并只给临时目录有限写权限。
容器不获得模型密钥、仓库认证或用户目录。
worker记录源码哈希、实参、退出码、stdout、stderr和请求输出是否存在；
独立分析器从序列化前的记录集合及精确算术计算参照。
这里测的是选定CLI的子契约和边界，不是整个自然包的安全或领域质量认证。

\subsection{外部SkillTV-Bench与中断后的分析范围}
按官方SkillTV-Bench评价划分固定全部203条轨迹、14个源任务，仓库提交为
\texttt{0774768c}\citep{S01}。
不使用其478条演化轨迹调参，也不按标签抽样。
标签和reward不进入模型输入，来源发现时查看过一个样例标签及公开类别总数，所以不是作者全程盲测。
推理只传原任务、轨迹白名单字段、技能正文及可选官方产物摘要。
三臂分别为普通直接评审、相同材料的证据规则、另加只读程序化产物事实。
共同输出成功概率、二元判断、短理由和引用定位；这是受限上下文的DeepSeek适配，不是原论文Harbor智能体或模型的完整复现。

轨迹最多48,000字符，技能合计12,000字符，原任务、官方摘要及新增产物事实分别最多12,000字符；
截断保留首尾并显式标记，不能把遗漏上下文当作原执行缺失。
程序只解析已有OBJ/JSON结构，不运行候选代码，不把OBJ语法事实当成渲染或任务正确证明。
公开快照本身可能不完整，未提供文件不等于原智能体没有生成文件。

计划609次，最多650次尝试，六并发，不自动重试，费用上限10美元。
使用官方 \texttt{deepseek-flash}、关闭思考、温度零、输出上限768 token；2026-11-13核对官方定价，
以每百万输入0.3美元、输出1.2美元的峰时单价保守核算。
第259次尝试返回402余额不足，守卫停止后续调度。
实际258个HTTP200响应，255个通过严格结构检查；第87个案例只启动了一臂失败请求，
完整配对为前86个排序案例，覆盖六个源任务、四个领域，其中24个官方通过、62个失败。
这不是预定203例的完成结果，也不是随机代表性子样本。

主指标为平衡准确率，结构无效响应按主分类错误计；另报准确率、错误放行、Brier与0.2/0.8置信处置的覆盖。
Brier中无效响应的概率固定为0.5并单列这种口径，不事后修补输出。
按实际六个任务簇重采样10,000次，区间只作条件描述。
事后检查显示，这86个案例没有官方产物摘要或可读产物文件，且66个轨迹被截断，
所以第三臂实际上主要增加了快照范围提示，\emph{尚未测试有产物可查时的增强收益}。
官方标签在预测和中断记录冻结后才用于评分；公开基准潜在训练污染未被排除。

\subsection{自然工具联合验收的探索性回放}
在已观察355次CLI结果后，另做明确的事后联合分析，目标为同输入的行数、原工具约定缺失统计和请求阈值的IQR计数全部正确。
沿用Iris/Wine开发、WDBC/Banknote来源留出，每边24个不同输入；两个阈值形成每边48个组合，显式回归例另计。
比较固定原始流程、原候选最短路径、开发失败过滤后的原候选路径，以及加入两处修复副本后的最短/证据路径。
路径复用已实际观察的只读CLI输出，不新增执行，不把回放次数计作新的独立任务。
该探索没有被注册为确认性网络终点，也没有用回放结果再修改补丁。
